\documentclass[11pt, a4paper]{article}
\usepackage{graphicx} % Required for inserting images

% Configuración de idioma y codificación
\usepackage[T1]{fontenc}
\usepackage[utf8]{inputenc}
% Mejorar compatibilidad con captions y hooks
\usepackage{caption}
\usepackage[spanish]{babel}
\usepackage{csquotes}

% Configuración de bibliografía
\usepackage[backend=biber, style=apa]{biblatex}
\addbibresource{referencias.bib}

% Paquetes para matemáticas e imágenes
\usepackage{amsmath}

\title{Documento LLMs, VLLMs}   
\author{Facundo Gaviola}
\date{\today}

\begin{document}

\maketitle


Este documento explora la evolución técnica en el campo del Procesamiento de Lenguaje Natural (NLP), analizando la transición desde las arquitecturas secuenciales basadas
en Redes Neuronales Recurrentes (RNN) y Long Short-Term Memory (LSTM) hasta el cambio de paradigma introducido por los Transformers y posteriores LLMs y su evolucion 
hacias los VLLMs.


\section{¿Qué es el Procesamiento de Lenguaje Natural (NLP)?}

El Procesamiento de Lenguaje Natural (NLP, por sus siglas en inglés) es una disciplina fundamental dentro de la Inteligencia Artificial (IA) que ha evolucionado rápidamente
en los últimos años. Su objetivo central es capacitar a las máquinas para comprender, interpretar y generar lenguaje humano de manera eficiente.

Según Atkinson-Abutridy, el NLP es el corazón de las aplicaciones innovadoras actuales, como los sistemas conversacionales (chatbots). La esencia de los avances modernos en
NLP radica en la capacidad de capturar patrones complejos y aprender representaciones contextuales del lenguaje, superando las limitaciones de los antiguos sistemas basados
en reglas fijas.

En la era actual, el NLP se centra predominantemente en el uso de \textbf{Grandes Modelos de Lenguaje (LLMs)}, los cuales son redes neuronales con miles de millones de
parámetros entrenados con cantidades masivas de texto no etiquetado mediante aprendizaje autosupervisado.

\subsection{Utilidades y Aplicaciones del NLP}

El libro destaca que los modelos de lenguaje modernos permiten a las computadoras realizar una amplia gama de tareas. Las aplicaciones principales descritas incluyen:

\begin{enumerate}
    \item \textbf{Sistemas Conversacionales y Chatbots:} Como ChatGPT o BARD, capaces de mantener diálogos coherentes, responder preguntas de seguimiento y recordar detalles 
    previos de la conversación.
    \item \textbf{Clasificación de Sentimientos:} Análisis de opiniones para determinar la polaridad emocional (positiva, negativa o neutra) en textos provenientes de redes 
    sociales o reseñas de productos.
    \item \textbf{Generación de Resúmenes:} Capacidad para sintetizar documentos largos en versiones concisas manteniendo la información clave.
    \item \textbf{Traducción Automática:} Traducción de texto entre diferentes idiomas con alta precisión contextual.
    \item \textbf{Búsqueda Semántica:} A diferencia de la búsqueda por palabras clave, los modelos de NLP modernos permiten buscar basándose en el significado y la intención 
    detrás de la consulta del usuario.
    \item \textbf{Generación de Código:} Herramientas que traducen descripciones en lenguaje natural a código de programación funcional (e.g., Python, SQL).
    \item \textbf{Razonamiento y QA (Question Answering):} Sistemas capaces de responder preguntas abiertas o cerradas basándose en un contexto dado o en conocimiento general.
\end{enumerate}

\subsection{Tipos de Modelos en NLP}

La evolución de los modelos en NLP ha pasado desde enfoques estadísticos hasta redes neuronales profundas y arquitecturas basadas en Transformers. A grandes rasgos, los modelos
se pueden caracterizar por su objetivo (generativos vs. discriminativos), por su arquitectura (encoders, decoders, o ambos) y por el tipo de tarea para la cual se optimizan 
(clasificación, generación, traducción, etc.).

Más allá de la distinción teórica, es útil considerar tres familias prácticas:
\begin{enumerate}
\item \textbf{Modelos basados en codificador (encoder-only):} están optimizados para comprender y representar texto. Estos modelos producen una representación contextual para
cada token y se usan principalmente en tareas de comprensión como clasificación, extracción de información y respuesta a preguntas. BERT y sus derivaciones son ejemplos clásicos.
\item \textbf{Modelos autorregresivos (decoder-only):} se entrenan para modelar la probabilidad condicional del siguiente token dado el contexto previo, lo que los hace especialmente
buenos para generación de texto y completado. GPT y sus sucesores encajan en este grupo.
\item \textbf{Modelos encoder-decoder (seq2seq):} combinan un encoder que resume la entrada y un decoder que genera la salida; son la elección natural para tareas de traducción 
automática, summarización y tareas de transformación de texto a texto. T5 y la mayoría de los sistemas de traducción neuronal usan este enfoque.
\end{enumerate}
Cada familia tiene ventajas y limitaciones: los encoders capturan contexto bidireccional y son fuertes en comprensión, los decoders autorregresivos resultan flexibles para generación
libre y continuaciones, y los encoder-decoder equilibran comprensión y generación para transformar una pieza de texto en otra.

\subsubsection{ Modelos Generativos vs. Discriminativos}
Es fundamental distinguir entre estas dos categorías:
\begin{itemize}
    \item \textbf{Modelos Discriminativos:} Se centran en clasificar datos o predecir una etiqueta $Y$ dada una entrada $X$ (calculan la probabilidad condicional $P(Y|X)$). 
    Ejemplo: Decidir si un correo es spam o no.
    \item \textbf{Modelos Generativos:} Aprenden la distribución de los datos y cómo se generan (calculan la probabilidad conjunta $P(X,Y)$). Tienen la capacidad de crear 
    nuevas instancias de datos (texto, imágenes) que se asemejan a los datos reales. Los LLMs modernos son fundamentalmente generativos.
\end{itemize}

\section{Redes Neuronales Recurrentes y LSTMs}
\subsection{El Procesamiento Secuencial}
Antes de los Transformers, el estándar para procesar secuencias eran las RNNs.
Las Redes Neuronales Recurrentes (RNN) se establecieron como el estado del arte para el modelado de secuencias. A diferencia de las redes tradicionales, las RNN procesan la información
de manera secuencial; es decir, procesan una palabra a la vez. Esto les permite mantener una "memoria" de los eventos pasados al establecer dependencias entre los puntos de datos 
\parencite{atkinson2025}. 

El desarrollo de las RNNs se fundamenta en el trabajo pionero de \textcite{rumelhart1986}, quienes popularizaron el algoritmo de \textbf{back-propagation}, esencial para el entrenamiento 
de redes profundas.\textbf{[EXPANDIR BACKPROPAGATION CUANDO TENGAS EL PAPER]} Posteriormente, \textcite{elman1990} introdujo las ``Redes de Elman'', que formalizaron el uso de capas de contexto para procesar secuencias temporales, permitiendo
a la red encontrar estructura en el tiempo. 

Sin embargo, las RNN enfrentan una limitación crítica conocida como el problema del \textit{\textbf{desvanecimiento del gradiente}} (vanishing gradient). A medida que la secuencia de entrada 
se alarga, los gradientes utilizados para actualizar los pesos durante el entrenamiento se vuelven extremadamente pequeños, impidiendo que la red aprenda patrones de dependencias a 
largo plazo. En términos prácticos, las RNN tienden a olvidar la información de secuencias de datos largas \parencite{atkinson2025}.
Para mitigar las limitaciones de las RNN simples, se introdujeron las redes LSTM. Estas arquitecturas poseen bloques de celdas más complejos que les permiten retener más entradas pasadas.
 Una LSTM gestiona el flujo de información a través de tres compuertas principales:
\begin{enumerate}
    \item \textbf{Compuerta de Olvido (Forget Gate):} Decide qué información del estado de la celda anterior ($C_{t-1}$) se debe descartar. Matemáticamente, evalúa la entrada actual y el estado oculto previo para generar un valor entre 0 (olvidar completamente) y 1 (mantener completamente).
    
    \item \textbf{Compuerta de Entrada (Input Gate):} Determina qué nueva información se almacenará en el estado de la celda. Utiliza una capa sigmoide para decidir qué valores actualizar y una capa `tanh` para crear un vector de nuevos candidatos.
    
    \item \textbf{Compuerta de Salida (Output Gate):} Controla qué partes del estado de la celda ($C_t$) se enviarán como salida ($h_t$). Esta salida es una versión filtrada del estado interno.
\end{enumerate}

A pesar de sus mejoras en la gestión del contexto, las LSTM y las RNN comparten una debilidad estructural inherente: su naturaleza secuencial impide la paralelización dentro de los ejemplos
de entrenamiento. El cálculo del estado oculto $h_t$ depende del estado anterior $h_{t-1}$, lo que hace que el tiempo de entrenamiento sea prohibitivo para secuencias muy largas y limite el 
uso eficiente del hardware moderno como las GPUs \parencite{vaswani2017}.

\section{La Revolución de los Transformers}

En 2017, Vaswani et al. propusieron una nueva arquitectura de red simple, el \textbf{Transformer}, basada únicamente en mecanismos de atención, prescindiendo por completo de la recurrencia 
y las convoluciones \parencite{vaswani2017}.

\subsection{Superando la Secuencialidad}
A diferencia de las RNN y LSTM, los Transformers procesan los datos en paralelo. Esto significa que leen todos los tokens de entrada a la vez en lugar de uno por uno. Esta capacidad de 
procesamiento no secuencial permite a los modelos realizar tareas de NLP mucho más rápido que sus predecesores \parencite{nvidia2023}.

Sin embargo, al eliminar la recurrencia, el modelo pierde la noción del orden de las palabras (por ejemplo, distinguir entre "El perro mordió al hombre" y "El hombre mordió al perro"). Para
solucionar esto, el Transformer inyecta información sobre la posición relativa o absoluta de los tokens en la secuencia mediante \textit{Positional Encodings} (codificaciones posicionales) 
\parencite{vaswani2017}.

\subsection{Mecanismo de Atención (Self-Attention)}
El núcleo de esta arquitectura es el mecanismo de "Auto-atención" (\textit{Self-Attention}). Este mecanismo permite al modelo relacionar diferentes posiciones de una sola secuencia para
calcular una representación de la misma. Matemáticamente, la función de atención se describe como el mapeo de una consulta ($Q$) y un conjunto de pares clave-valor ($K, V$) a una salida:

\begin{equation}
    Attention(Q, K, V) = \text{softmax}\left(\frac{QK^T}{\sqrt{d_k}}\right)V
\end{equation}

Donde $d_k$ es la dimensión de las claves. Este cálculo permite al modelo ponderar la importancia de cada palabra en relación con todas las demás palabras de la oración, independientemente de
su distancia física en la secuencia, resolviendo así el problema de las dependencias a largo plazo que afectaba a las RNN \parencite{vaswani2017}.

Esta arquitectura permitió el desarrollo de Modelos de Lenguaje Grandes (LLMs) pre-entrenados como BERT y GPT, los cuales utilizan cantidades masivas de datos para capturar patrones complejos
del lenguaje \parencite{atkinson2025}.


\printbibliography

\end{document}
